\documentclass[11pt,a4paper]{article}

\usepackage[T1]{fontenc}
\usepackage[utf8]{inputenc}
\usepackage{lmodern}
\usepackage[a4paper,margin=2.2cm]{geometry}
\usepackage{setspace}
\usepackage{parskip}
\usepackage{enumitem}
\usepackage{hyperref}
\usepackage[nameinlink,noabbrev]{cleveref}

\setstretch{1.05}
\setlist[itemize]{topsep=3pt,itemsep=2pt}
\setlist[enumerate]{topsep=3pt,itemsep=2pt}

\title{Antihypertensives\\Structured Rewrite of Peck and Harris, Chapter 17}
\author{}
\date{Source: \textit{Pharmacology for Anaesthesia and Intensive Care}, 5th ed., Chapter 17 (AntiHypertensives)}

\begin{document}

\maketitle
\tableofcontents
\newpage

\section{Core framework}

This chapter groups antihypertensive drugs into:
\begin{enumerate}
    \item drugs acting on the renin--angiotensin--aldosterone system (RAAS),
    \item adrenergic neurone-blocking drugs,
    \item centrally acting antihypertensives.
\end{enumerate}

Within this chapter, the principal drugs discussed are:
\begin{itemize}
    \item ACE inhibitors,
    \item angiotensin II receptor antagonists (ARAs),
    \item guanethidine,
    \item metirosine,
    \item methyldopa,
    \item clonidine.
\end{itemize}

Diuretics, adrenoceptor antagonists, and calcium channel antagonists are noted as important antihypertensive groups, but are addressed elsewhere rather than in detail here.

\section{Renin--angiotensin--aldosterone system}

\subsection{Physiological basis}

The juxtaglomerular apparatus consists of:
\begin{itemize}
    \item juxtaglomerular cells in the afferent arteriole, containing prorenin/renin,
    \item macula densa cells at the beginning of the distal convoluted tubule,
    \item agranular lacis cells between the afferent and efferent arterioles.
\end{itemize}

Renin release is stimulated by:
\begin{itemize}
    \item reduced renal perfusion,
    \item reduced sodium delivery to the macula densa,
    \item sympathetic stimulation via $\beta_1$-adrenoceptors.
\end{itemize}

Renin converts angiotensinogen to angiotensin~I. Angiotensin-converting enzyme (ACE) converts angiotensin~I to angiotensin~II and also degrades bradykinin. Angiotensin~II is further metabolised to inactive products and angiotensin~III, which retains some biological activity.

\subsection{Actions of angiotensin II}

Angiotensin~II acts predominantly at AT$_1$ receptors and produces:
\begin{itemize}
    \item potent vasoconstriction,
    \item facilitation of sympathetic activity, including reduced noradrenaline reuptake,
    \item increased thirst,
    \item increased ADH and ACTH release,
    \item increased aldosterone secretion,
    \item suppression of renin release,
    \item reduced glomerular filtration rate.
\end{itemize}

These actions explain why RAAS-targeting drugs are useful in hypertension and heart failure.

\section{ACE inhibitors}

\subsection{Clinical role}

ACE inhibitors are used in:
\begin{itemize}
    \item all grades of heart failure,
    \item post-myocardial infarction patients with left ventricular dysfunction,
    \item hypertension,
    \item diabetic nephropathy, particularly in insulin-dependent diabetes.
\end{itemize}

The chapter notes that hypertension may be less responsive to ACE inhibition in black patients, in whom concurrent diuretic therapy may be required.

A key anaesthetic point is that although most chronic drugs are continued peri-operatively, ACE inhibitors and ARAs are often omitted because they increase the risk of peri-operative hypotension.

\subsection{Kinetic classification}

ACE inhibitors are grouped kinetically into:
\begin{itemize}
    \item \textbf{Group 1:} captopril --- active drug, metabolised to active metabolites,
    \item \textbf{Group 2:} enalapril and ramipril --- prodrugs activated by hepatic metabolism to the diacid form,
    \item \textbf{Group 3:} lisinopril --- active drug, not metabolised, excreted unchanged in urine.
\end{itemize}

\section{Ramipril}

\subsection{Presentation}

Ramipril is available as tablets of 1.25--10~mg. Treatment is started at 1.25~mg and titrated according to indication.

\subsection{Mechanism of action}

Ramipril is converted by liver esterases to ramiprilat, a competitive ACE inhibitor. ACE inhibition reduces formation of angiotensin~II and thereby reduces its cardiovascular and endocrine effects. The chapter states that afterload falls more than preload.

\subsection{Pharmacological effects}

\subsubsection{Cardiovascular effects}

Ramipril:
\begin{itemize}
    \item decreases systemic vascular resistance,
    \item lowers arterial blood pressure,
    \item may increase cardiac output in heart failure because afterload is reduced,
    \item usually does not change heart rate substantially, though a rise may occur,
    \item preserves baroreflex function,
    \item may produce transient hypotension when treatment is initiated.
\end{itemize}

\subsubsection{Renal effects}

Angiotensin~II normally helps preserve glomerular filtration in low-perfusion states by constricting the efferent arteriole. ACE inhibition removes this compensatory mechanism. Consequently:
\begin{itemize}
    \item renal perfusion pressure may fall,
    \item renal failure may develop in susceptible patients,
    \item bilateral renal artery stenosis, or unilateral stenosis in a solitary functioning kidney, is a contraindication.
\end{itemize}

Where renal perfusion is otherwise adequate, renal vasodilation and natriuresis may occur.

\subsubsection{Metabolic effects}

By reducing aldosterone secretion, ramipril can cause:
\begin{itemize}
    \item hyperkalaemia,
    \item increases in urea and creatinine, especially in pre-existing renal impairment.
\end{itemize}

\subsubsection{Adverse effects and interactions}

Important adverse effects and interactions include:
\begin{itemize}
    \item hyperkalaemia, especially with potassium-sparing diuretics,
    \item unexplained hypoglycaemia in both type~1 and type~2 diabetes,
    \item reduced antihypertensive efficacy and increased renal risk with NSAIDs,
    \item dry cough due to reduced bradykinin breakdown,
    \item rare angioedema,
    \item rare agranulocytosis and thrombocytopenia,
    \item loss of taste, rash, pruritus, fever, and aphthous ulceration.
\end{itemize}

\subsection{Pharmacokinetics}

Ramipril:
\begin{itemize}
    \item is a prodrug,
    \item is rapidly absorbed orally,
    \item yields ramiprilat with oral bioavailability of about 50\%,
    \item has a half-life of about 15~hours at high doses, longer at low doses because ACE binding is saturable,
    \item is excreted mainly by the kidney, so elimination is prolonged in renal failure.
\end{itemize}

\section{Enalapril}

Enalapril is also a prodrug and is converted to enalaprilat in the liver and kidney. It may be administered orally as enalapril or intravenously as enalaprilat.

Key pharmacokinetic points:
\begin{itemize}
    \item elimination half-life 4--8~hours,
    \item during prolonged therapy this may rise to about 11~hours,
    \item duration of action is approximately 20~hours.
\end{itemize}

\section{Angiotensin II receptor antagonists}

\subsection{General principle}

Angiotensin II receptor antagonists (ARAs) block the action of angiotensin~II directly at the receptor rather than preventing its formation.

\section{Losartan}

\subsection{Presentation and indications}

Losartan is available as 25--50~mg tablets and also in combination with hydrochlorothiazide. It is used in hypertension, especially when ACE inhibitor-induced cough limits therapy.

\subsection{Mechanism of action}

Losartan is a selective AT$_1$ receptor antagonist. Blockade of AT$_1$ receptors removes the negative feedback of angiotensin~II on renin secretion. As a result, renin and angiotensin~II concentrations rise, but the increased angiotensin~II has little effect because AT$_1$ receptors remain blocked.

\subsection{Effects}

Its haemodynamic effects are broadly similar to those of ACE inhibitors. However, because ACE is not inhibited:
\begin{itemize}
    \item bradykinin metabolism remains intact,
    \item dry cough is not a characteristic adverse effect.
\end{itemize}

\subsection{Pharmacokinetics}

Losartan:
\begin{itemize}
    \item is a substituted imidazole compound,
    \item behaves effectively as a prodrug because its active metabolite is much more potent,
    \item is well absorbed orally,
    \item undergoes extensive first-pass metabolism to an active carboxylic acid metabolite and inactive metabolites,
    \item has oral bioavailability of about 30\%,
    \item is 99\% protein bound,
    \item has a parent drug half-life of about 2~hours,
    \item has an active metabolite half-life of about 7~hours,
    \item has less than 10\% urinary excretion as unchanged active drug,
    \item excretes inactive metabolites in bile and urine.
\end{itemize}

\subsection{Contraindications}

Contraindications stated in the source include:
\begin{itemize}
    \item bilateral renal artery stenosis,
    \item pregnancy.
\end{itemize}

\subsection{Why ARAs may be preferred over ACE inhibitors}

The source gives three advantages:
\begin{enumerate}
    \item AT$_1$ blockade is a more specific method of preventing the harmful effects of angiotensin~II,
    \item AT$_2$ receptors remain unblocked and may have cardioprotective effects,
    \item cough and angioedema are less common, which may improve compliance.
\end{enumerate}

\section{Adrenergic neurone-blocking drugs}

\subsection{Physiological basis}

Noradrenaline is synthesised from tyrosine within adrenergic nerve terminals, stored in vesicles, and released into the synaptic cleft.

Termination of its action occurs mainly by:
\begin{itemize}
    \item \textbf{Uptake 1:} reuptake into the nerve terminal by a high-affinity transporter,
    \item \textbf{Uptake 2:} diffusion away into the circulation and extraneuronal uptake, especially in the liver, with metabolism by catechol-\emph{O}-methyltransferase.
\end{itemize}

\section{Guanethidine}

\subsection{Presentation and uses}

Guanethidine is supplied as tablets and as a colourless injection of 10~mg/mL. It was formerly used as an antihypertensive, but is now mainly used for sympathetically mediated chronic pain.

The chapter gives:
\begin{itemize}
    \item antihypertensive dosing: initially 10~mg/day, then 30--50~mg/day,
    \item intravenous regional block dose for chronic pain: 20~mg.
\end{itemize}

\subsection{Mechanism of action}

Guanethidine enters adrenergic neurones via uptake~1.

After intravenous administration, a triphasic response may occur:
\begin{enumerate}
    \item initial hypotension from direct arteriolar vasodilation,
    \item transient hypertension due to displacement of noradrenaline,
    \item sustained hypotension because release of remaining noradrenaline is prevented.
\end{enumerate}

This pattern is not seen with oral dosing because onset is slower. Catecholamine secretion from the adrenal medulla is not affected.

\subsection{Effects and adverse effects}

Major effects include:
\begin{itemize}
    \item hypotension,
    \item marked postural hypotension,
    \item fluid retention and oedema,
    \item diarrhoea,
    \item failure of ejaculation.
\end{itemize}

Its action is antagonised by drugs that block uptake~1, such as:
\begin{itemize}
    \item tricyclic antidepressants,
    \item cocaine.
\end{itemize}

The source also notes that chronic administration causes adrenoceptor up-regulation, which makes patients unusually sensitive to direct-acting sympathomimetics.

\subsection{Pharmacokinetics}

Guanethidine:
\begin{itemize}
    \item is variably and incompletely absorbed orally,
    \item undergoes first-pass metabolism, giving oral bioavailability of about 50\%,
    \item is not protein bound,
    \item does not cross the blood--brain barrier,
    \item has an elimination half-life of several days,
    \item is eliminated by hepatic metabolism and renal excretion of unchanged drug and metabolites.
\end{itemize}

\section{Metirosine}

Metirosine is a competitive inhibitor of tyrosine hydroxylase and therefore reduces catecholamine synthesis. The source states that it should only be used for hypertension associated with phaeochromocytoma.

Adverse effects listed in the chapter are:
\begin{itemize}
    \item severe diarrhoea,
    \item sedation,
    \item extrapyramidal effects,
    \item hypersensitivity reactions.
\end{itemize}

\section{Centrally acting antihypertensives}

These drugs lower blood pressure primarily by reducing central sympathetic outflow.

\section{Methyldopa}

\subsection{Presentation and use}

Methyldopa is available as:
\begin{itemize}
    \item film-coated tablets of 125--500~mg,
    \item oral suspension 50~mg/mL,
    \item intravenous preparation 50~mg/mL, although this is rarely used.
\end{itemize}

Typical dosing is 250~mg three times daily, increasing up to 3~g/day. It is especially important in hypertension associated with pregnancy.

\subsection{Mechanism of action}

Methyldopa crosses the blood--brain barrier and is converted to $\alpha$-methyl-noradrenaline. This metabolite is a potent $\alpha_2$-agonist with limited $\alpha_1$ activity.

By stimulating presynaptic $\alpha_2$ receptors in the nucleus tractus solitarius, it reduces further noradrenaline release and decreases centrally mediated sympathetic tone.

\subsection{Effects}

\subsubsection{Cardiovascular effects}

Methyldopa:
\begin{itemize}
    \item reduces systemic vascular resistance,
    \item lowers arterial pressure,
    \item may cause postural hypotension,
    \item leaves cardiac output unchanged despite relative bradycardia,
    \item may produce rebound hypertension if stopped abruptly, although this is less common than with clonidine.
\end{itemize}

\subsubsection{Central nervous system effects}

Methyldopa commonly causes sedation. Less common CNS effects include:
\begin{itemize}
    \item dizziness,
    \item depression,
    \item nightmares,
    \item nausea.
\end{itemize}

It also reduces the MAC of volatile anaesthetic agents.

\subsubsection{Haematological and immune effects}

Important effects include:
\begin{itemize}
    \item positive direct Coombs' test in 10--20\%,
    \item rare thrombocytopenia and leucopenia,
    \item autoimmune haemolytic anaemia,
    \item eosinophilia with fever early in treatment,
    \item hypersensitivity myocarditis.
\end{itemize}

\subsubsection{Renal and hepatic effects}

Additional effects include:
\begin{itemize}
    \item urine darkening on exposure to air,
    \item deterioration of liver function during prolonged therapy,
    \item rare fatal hepatic necrosis.
\end{itemize}

\subsubsection{Miscellaneous effects}

Methyldopa:
\begin{itemize}
    \item may falsely elevate urinary catecholamine assays because it fluoresces at similar wavelengths,
    \item does not affect VMA assays,
    \item may cause constipation,
    \item may cause gynaecomastia.
\end{itemize}

\subsection{Pharmacokinetics}

Methyldopa:
\begin{itemize}
    \item is absorbed erratically from the gut,
    \item has very variable oral bioavailability,
    \item has a slow onset,
    \item undergoes first-pass hepatic metabolism to the \emph{O}-sulfate,
    \item is less than 20\% protein bound,
    \item is excreted about 50\% unchanged in urine.
\end{itemize}

\section{Clonidine}

\subsection{Pharmacological identity}

Clonidine is an $\alpha$-agonist with approximately 200-fold greater affinity for $\alpha_2$ than for $\alpha_1$ receptors. Some studies identify it as a partial agonist.

\subsection{Presentation and uses}

Clonidine is available as:
\begin{itemize}
    \item tablets of 25--300~$\mu$g,
    \item injection 150~$\mu$g/mL,
    \item transdermal patches, although therapeutic levels take about 48~hours to develop.
\end{itemize}

Uses listed in the source are:
\begin{itemize}
    \item hypertension,
    \item acute and chronic pain,
    \item suppression of opioid withdrawal symptoms,
    \item augmentation of sedation during ventilation in critical illness.
\end{itemize}

\subsection{Mechanism of action}

The main useful actions of clonidine arise from:
\begin{itemize}
    \item $\alpha_2$ stimulation in the lateral reticular nucleus, reducing central sympathetic outflow,
    \item $\alpha_2$ stimulation in the spinal cord, enhancing endogenous opioid release and modulating descending noradrenergic pathways involved in nociception.
\end{itemize}

The reduction in volatile anaesthetic MAC appears to be mediated by central postsynaptic $\alpha_2$ receptors. $\alpha_2$ receptor activation is coupled to G$_i$, reduces intracellular cAMP, and also activates potassium channels.

\subsection{Effects}

\subsubsection{Cardiovascular effects}

Clonidine:
\begin{itemize}
    \item may initially increase blood pressure after intravenous administration because of peripheral $\alpha_1$ stimulation,
    \item is then followed by a longer hypotensive phase,
    \item usually maintains cardiac output despite bradycardia,
    \item prolongs the PR interval,
    \item depresses AV nodal conduction,
    \item sensitises baroreceptor reflexes,
    \item stabilises haemodynamic responses to peri-operative stimuli,
    \item may cause rebound hypertension if abruptly stopped, especially at doses above 1.2~mg/day,
    \item may have rebound worsened by concurrent $\beta$-blockade because vasoconstriction becomes unopposed,
    \item shows a ceiling effect at higher doses because $\alpha_1$ stimulation becomes limiting.
\end{itemize}

\subsubsection{Central nervous system effects}

Clonidine:
\begin{itemize}
    \item causes sedation,
    \item may reduce volatile anaesthetic MAC by up to 50\%,
    \item is anxiolytic at low doses and anxiogenic at high doses.
\end{itemize}

\subsubsection{Analgesic effects}

Clonidine:
\begin{itemize}
    \item may be given intrathecally or epidurally,
    \item provides prolonged analgesia,
    \item does not cause respiratory depression,
    \item appears synergistic with opioids,
    \item does not itself produce sensory or motor block,
    \item may reduce postoperative opioid requirements when given intravenously.
\end{itemize}

\subsubsection{Renal effects}

It may cause diuresis, for which inhibition of ADH release is one suggested mechanism.

\subsubsection{Respiratory effects}

Unlike opioids, clonidine does not produce significant respiratory depression.

\subsubsection{Endocrine effects}

Clonidine:
\begin{itemize}
    \item blunts the stress response to surgery,
    \item inhibits insulin release, although this rarely causes important hyperglycaemia.
\end{itemize}

\subsubsection{Haematological effects}

Although platelets possess $\alpha_2$ receptors, therapeutic doses do not promote platelet aggregation. Sympatholysis blocks adrenaline-induced platelet aggregation.

\subsection{Pharmacokinetics}

Clonidine:
\begin{itemize}
    \item is rapidly and almost completely absorbed orally,
    \item has oral bioavailability close to 100\%,
    \item is 20\% protein bound,
    \item has a volume of distribution of about 2~L/kg,
    \item has an elimination half-life of 9--18~hours,
    \item undergoes about 50\% hepatic metabolism to inactive metabolites,
    \item is otherwise excreted unchanged in urine,
    \item requires dose reduction in renal impairment.
\end{itemize}

\section{Atipamezole}

Atipamezole is a selective $\alpha_2$ antagonist that crosses the blood--brain barrier and reverses the sedation and analgesia produced by clonidine and dexmedetomidine. Its elimination half-life is about 2~hours, similar to dexmedetomidine.

\section{Anaesthesia-focused takeaways}

The following points are especially relevant for anaesthetic practice:
\begin{itemize}
    \item ACE inhibitors and ARAs are often withheld peri-operatively because they increase the risk of hypotension.
    \item ACE inhibitors and ARAs may precipitate renal failure in renal artery stenosis because they remove angiotensin~II support of glomerular filtration.
    \item Methyldopa and clonidine both reduce volatile anaesthetic MAC.
    \item Clonidine has important peri-operative roles beyond antihypertensive therapy, including sedation, analgesia, opioid sparing, blunting of haemodynamic responses, and attenuation of the surgical stress response.
    \item Long-term guanethidine therapy leads to adrenoceptor up-regulation, causing increased sensitivity to direct-acting sympathomimetics.
\end{itemize}

\section{Condensed exam summary}

\subsection{ACE inhibitors}
\begin{itemize}
    \item inhibit formation of angiotensin~II,
    \item reduce systemic vascular resistance and afterload,
    \item are used in heart failure, post-myocardial infarction left ventricular dysfunction, hypertension, and diabetic nephropathy,
    \item can cause hypotension, renal impairment, hyperkalaemia, cough, and angioedema,
    \item are contraindicated in significant renal artery stenosis,
    \item are often omitted peri-operatively.
\end{itemize}

\subsection{Angiotensin II receptor antagonists}
\begin{itemize}
    \item block AT$_1$ receptors directly,
    \item have effects broadly similar to ACE inhibitors,
    \item do not inhibit bradykinin breakdown, so cough is much less common,
    \item are contraindicated in renal artery stenosis and pregnancy.
\end{itemize}

\subsection{Guanethidine}
\begin{itemize}
    \item is an adrenergic neurone blocker entering via uptake~1,
    \item now has limited antihypertensive use,
    \item causes postural hypotension,
    \item interacts with drugs that inhibit uptake~1,
    \item increases sensitivity to direct-acting sympathomimetics after chronic use.
\end{itemize}

\subsection{Metirosine}
\begin{itemize}
    \item inhibits tyrosine hydroxylase,
    \item is reserved for hypertension associated with phaeochromocytoma.
\end{itemize}

\subsection{Methyldopa}
\begin{itemize}
    \item is converted centrally to an $\alpha_2$-agonist,
    \item is important in hypertension during pregnancy,
    \item causes sedation,
    \item may cause a positive Coombs' test, haemolysis, and hepatic toxicity.
\end{itemize}

\subsection{Clonidine}
\begin{itemize}
    \item is a potent $\alpha_2$-agonist,
    \item acts as an antihypertensive, sedative, and analgesic adjunct,
    \item reduces volatile anaesthetic MAC,
    \item blunts stress responses,
    \item does not cause major respiratory depression,
    \item may cause rebound hypertension if stopped abruptly.
\end{itemize}

\end{document}
